\documentclass{bmstu}

\begin{document}

\makereporttitle
{Информатика и системы управления (ИУ)}
{Программное обеспечение ЭВМ и информационные технологии (ИУ7)}
{\textbf{4}}
{Проектирование экспертных систем}
{Процедура резолюции с унификацией}
{}
{ИУ7-32М}
{К.Э. Ковалец}
{З.Н. Русакова}


\setcounter{page}{2}
\renewcommand{\contentsname}{Содержание} 
\tableofcontents


\chapter{Основная часть}

\section{Унификация}

Унификация --- это процесс замены предметных переменных аргумента атома другими предметными переменными или константами или функциональными символами. Подстановка называется унификацией, если атомы совпадают.

Унификация выполняется:
\begin{itemize}
    \item если термы --- константы, то они унифицируются, если совпадают;
    \item если в первом атоме терм --- переменная, во втором --- константа, то они унифицируются, и переменная получает значение константы;
    \item если терм в первом атоме --- переменная и во втором --- переменная, то они унифицируются и становятся связанными, можно использовать любое имя.
\end{itemize}

Унифицируются между собой термы, стоящие на одинаковых местах в контрарных атомах.

\section{Резолюция}

Резольвента --- это объединение двух дизъюнктов за исключением контрарной пары.
Метод формирования резольвенты заключается в следующем:
\begin{itemize}
    \item сначала для двух атомов формируется список подстановок с помощью процедуры унификации;
    \item затем эти значения распространяются в атомы первого и второго дизъюнктов,
    за исключением контрарной пары.
\end{itemize}

Множество подстановок при унификации атомов и дизъюнктов нужно выполнять последовательно, просматривая каждый раз только одну переменную.

\section{Алгоритм вывода по методу резолюции}

Необходимо создать базу данных, которая будет представлять собой список дизъюнктов. 

Алгоритм состоит из следующих шагов:
\begin{itemize}
    \item принять отрицание заключения (того, что хотим доказать);
    \item привести все формулы аксиом и отрицание цели к КНФ;
    \item получить список дизьюнктов;
    \item если существует пара дизьюнктов, содержащих контрарные атомы, то эти дизьюнкты объединить с удалением контрарной пары;
    \item если на каком-то шаге получим два дизьюнкта, содержащих одинаковые атомы с разными знаками, то резольвентой будет пустой дизьюнкт или ложь;
    \item иначе продолжить формирование резольвенты (перебор на множестве).
\end{itemize}

\section{Пример работы программы}

Результаты работы программы приведены в листинге \ref{lst:result}.

\mylisting[text]{result.txt}{firstline=1,lastline=82}{Результат работы программы}{result}{}

\chapter{Текст программы}

В листингах \ref{lst:main}--\ref{lst:utils} представлен код программы.

\mylisting[python]{main.py}{firstline=1,lastline=69}{Основной модуль программы}{main}{}

\mylisting[python]{items.py}{firstline=1,lastline=52}{Модуль с описанием классов термов, атомов и дизъюнктов}{items}{}

\mylisting[python]{unification.py}{firstline=1,lastline=143}{Модуль с реализацией класса унификации}{unification}{}

\mylisting[python]{utils.py}{firstline=1,lastline=27}{Модуль с реализацией функции парсера дизъюнкта}{utils}{}

\end{document}
